\PassOptionsToPackage{table}{xcolor}
\documentclass[10pt,aspectratio=169]{beamer}
\newcommand{\LectureNumber}{26}
\input{course-theme.tex}
\title{Programming environments and debugging}
\subtitle{CSC103 Programming Fundamentals / Lecture 26}
\author{Dr. Muhammad Farhan}
\institute{Associate Professor of Computer Science\\Department of Computer Science\\COMSATS University Islamabad\\Sahiwal Campus}
\date{Fall 2026}
\begin{document}
\begin{frame}\titlepage\end{frame}

\begin{frame}[fragile,t]{The problem for this lecture}
\begin{columns}[T,onlytextwidth]\begin{column}{0.60\textwidth}
Extract the demonstrated maximum calculation into max(int[] values). Specify empty-input behaviour and test negative-only data.
\end{column}\begin{column}{0.35\textwidth}\small
Document that the method requires a nonempty array.
\end{column}\end{columns}
\note{Introduce the target problem before explaining the tools. Revisit it in the guided solution later.\par Syllabus reading: Reference material. CLO-3.}
\end{frame}

\begin{frame}[fragile,t]{Learning objectives}
\begin{columns}[T,onlytextwidth]\begin{column}{0.60\textwidth}
\begin{itemize}
\item Use breakpoints and variable inspection
\item Distinguish debugging from refactoring
\item Refactor while preserving observable behaviour
\end{itemize}
\end{column}\begin{column}{0.35\textwidth}\small
CLO-3

A modern programming environment\par Debugger operations\par Refactoring\par A repeatable debugging process
\end{column}\end{columns}
\note{Explain how each objective supports the opening problem. The final questions check these objectives.\par Syllabus reading: Reference material. CLO-3.}
\end{frame}

\begin{frame}[fragile,t]{A modern programming environment}
\begin{itemize}
\item An IDE integrates editing, compilation, navigation and debugging.
\item Project settings select a JDK and source level.
\item A clean build reduces confusion from stale compiled files.
\end{itemize}
\vspace{1mm}\begin{block}{Example}
\begin{lstlisting}
Before a demo:
Check the active JDK.
Check the working directory.
Run the intended main class.
\end{lstlisting}
\end{block}
\note{Use the locally available IDE without requiring a paid product. An editor and terminal also suffice for the code in these decks.\par Syllabus reading: Reference material. CLO-3.}
\end{frame}

\begin{frame}[fragile,t]{Debugger operations}
\begin{itemize}
\item A breakpoint pauses before a selected statement executes.
\item Step over executes a call without entering it. Step into enters a call.
\item Inspect variables and stack frames to test a hypothesis.
\end{itemize}
\vspace{1mm}\begin{block}{Example}
\begin{lstlisting}
Hypothesis: the loop misses its final element.
Breakpoint: loop condition
Watch: i, values.length, total
\end{lstlisting}
\end{block}
\note{Show actual state before changing code. Explain that a debugger changes timing and is a diagnostic tool, not a substitute for understanding.\par Syllabus reading: Reference material. CLO-3.}
\end{frame}

\begin{frame}[fragile,t]{Refactoring}
\begin{itemize}
\item Refactoring improves internal structure while preserving intended observable behaviour.
\item Examples include renaming and extracting a method.
\item Run existing tests before and after the change.
\end{itemize}
\vspace{1mm}\begin{block}{Example}
\begin{lstlisting}
Before: repeated discount formula
After: one discount(price, rate) method
Same inputs should produce the same outputs.
\end{lstlisting}
\end{block}
\note{Fixing an incorrect requirement or adding a feature is a behaviour change. A refactoring can accompany a bug fix, but evaluate them separately.\par Syllabus reading: Reference material. CLO-3.}
\end{frame}

\begin{frame}[fragile,t]{A repeatable debugging process}
\begin{itemize}
\item Reproduce the failure with the smallest useful input.
\item Predict the faulty state, inspect it and make a focused correction.
\item Keep a regression case that demonstrates the original bug.
\end{itemize}
\vspace{1mm}\begin{block}{Example}
\begin{lstlisting}
Failure: max({-4,-2}) returns 0.
Cause: max starts at 0.
Correction: initialise from the first element.
\end{lstlisting}
\end{block}
\note{This distinguishes guessing from evidence. The method needs a nonempty-array contract before using its first element.\par Syllabus reading: Reference material. CLO-3.}
\end{frame}

\begin{frame}[fragile,t]{A minimal failing case guides a repair}
\begin{itemize}
\item The smallest example that violates the contract is easier to reason about.
\item For a maximum function initialised to zero, \{-2\} already proves the defect.
\item The repair needs an empty-input policy before using the first element.
\end{itemize}
\vspace{1mm}\begin{block}{Example}
\begin{lstlisting}
Failure: max({-2}) returns 0.
Required result: -2.
A positive-only test hides this defect.
\end{lstlisting}
\end{block}
\note{Explain the mechanism using the displayed example. The smallest example that violates the contract is easier to reason about. For a maximum function initialised to zero, \{-2\} already proves the defect. The repair needs an empty-input policy before using the first element.\par Syllabus reading: Reference material. CLO-3.}
\end{frame}

\begin{frame}[fragile,t]{Extraction should preserve the calculation}
\begin{itemize}
\item Choose the statements that implement one responsibility.
\item Replace their external inputs with parameters and their result with a return value.
\item Run the same examples before and after extraction to check behaviour.
\end{itemize}
\vspace{1mm}\begin{block}{Example}
\begin{lstlisting}
Before: loop embedded in main
After: main calls max(values)
The same array should yield the same maximum.
\end{lstlisting}
\end{block}
\note{Explain the mechanism using the displayed example. Choose the statements that implement one responsibility. Replace their external inputs with parameters and their result with a return value. Run the same examples before and after extraction to check behaviour.\par Syllabus reading: Reference material. CLO-3.}
\end{frame}


\begin{frame}[fragile,t]{Debugging is a hypothesis-checking cycle}
\begin{center}\begin{tikzpicture}
\node[box] (n0) at (0.0,0) {Reproduce};
\node[box] (n1) at (3.45,0) {Pause and inspect};
\node[box] (n2) at (6.9,0) {Explain cause};
\node[alt] (n3) at (10.350000000000001,0) {Fix and retest};
\draw[arr] (n0) -- (n1);
\draw[arr] (n1) -- (n2);
\draw[arr] (n2) -- (n3);
\draw[arr] (n3.south) -- ++(0,-.8) -| (n0.south);
\end{tikzpicture}\end{center}
\begin{block}{What to notice}Inspect the first state that differs from your prediction, then rerun the reproducing case.\end{block}
\note{Trace each labelled connection. Inspect the first state that differs from your prediction, then rerun the reproducing case.}
\end{frame}
\begin{frame}[fragile,t]{Key distinctions}
\small\renewcommand{\arraystretch}{1.5}
\rowcolors{2}{pale}{white}
\begin{tabularx}{\textwidth}{>{\raggedright\arraybackslash}X>{\raggedright\arraybackslash}X>{\raggedright\arraybackslash}X}
\rowcolor{navy}
\color{white}\bfseries Observation & \color{white}\bfseries Hypothesis & \color{white}\bfseries Debugger check\\
Negative-only case gives 0 & Bad initial candidate & Inspect maximum before loop\\
Last value ignored & Wrong upper bound & Watch final index\\
Total changes between calls & Shared state retained & Inspect field or accumulator\\
\end{tabularx}
\vspace{3mm}\footnotesize These distinctions determine which operation or control structure fits the problem.
\note{\par Syllabus reading: Reference material. CLO-3.}
\end{frame}

\begin{frame}[fragile,t]{First example: execution}
\begin{lstlisting}
int[] values = {-4, -2, -9};
int maximum = values[0];
for (int i = 1; i < values.length; i++) {
  if (values[i] > maximum) maximum = values[i];
}
System.out.println(maximum);
\end{lstlisting}
\note{This is the main body from Java\_Examples/Lecture26.java. The source file includes the complete class. Predict the result before the next slide.\par Syllabus reading: Reference material. CLO-3.}
\end{frame}

\begin{frame}[fragile,t]{First example: result}
\begin{columns}[T,onlytextwidth]\begin{column}{0.60\textwidth}
-2\par Initial maximum: -4. After comparing -2: -2.\par Comparing -9 leaves the maximum unchanged.
\end{column}\begin{column}{0.35\textwidth}\small
Refactoring

A repeatable debugging process
\end{column}\end{columns}
\note{Expected output and interpretation: -2
Initial maximum: -4. After comparing -2: -2.
Comparing -9 leaves the maximum unchanged.\par Syllabus reading: Reference material. CLO-3.}
\end{frame}

\begin{frame}[fragile,t]{Guided practice}
\begin{columns}[T,onlytextwidth]\begin{column}{0.60\textwidth}
Extract the demonstrated maximum calculation into max(int[] values). Specify empty-input behaviour and test negative-only data.
\end{column}\begin{column}{0.35\textwidth}\small
Attempt the solution before continuing.

Use the definitions and examples from this lecture.
\end{column}\end{columns}
\note{Give students 8 minutes to attempt the problem. The following slides provide a complete solution. Original solution guidance: Reject an empty array with IllegalArgumentException. Start from values[0], traverse from index 1 and return maximum. \{-4,-2,-9\} must give -2.\par Syllabus reading: Reference material. CLO-3.}
\end{frame}

\begin{frame}[fragile,t]{Solution strategy}
\begin{columns}[T,onlytextwidth]\begin{column}{0.60\textwidth}
\begin{itemize}
\item Document that the method requires a nonempty array.
\item Move the maximum calculation into a method returning int.
\item Call it for negative-only data and reject the empty case explicitly.
\end{itemize}
\end{column}\begin{column}{0.35\textwidth}\small
The next program uses fixed inputs so its result can be reproduced.
\end{column}\end{columns}
\note{Discuss why each step is needed. State the input assumptions before executing the program.\par Syllabus reading: Reference material. CLO-3.}
\end{frame}

\begin{frame}[fragile,t]{Complete solution (1/2)}
\begin{lstlisting}
public class Solution26 {
  public static void main(String[] args) throws Exception {
    System.out.println(max(new int[]{-4, -2, -9}));
    try {
      max(new int[]{});
    } catch (IllegalArgumentException ex) {
      System.out.println("Empty array rejected");
    }
  }
  static int max(int[] values) {
\end{lstlisting}
\note{Complete runnable file: Java\_Examples/Solution26.java. Inputs are fixed in the program.  \par Syllabus reading: Reference material. CLO-3.}
\end{frame}

\begin{frame}[fragile,t]{Complete solution (2/2)}
\begin{lstlisting}
    if (values == null || values.length == 0) {
      throw new IllegalArgumentException("nonempty array required");
    }
    int maximum = values[0];
    for (int i = 1; i < values.length; i++) {
      if (values[i] > maximum) maximum = values[i];
    }
    return maximum;
  }
}
\end{lstlisting}
\note{Complete runnable file: Java\_Examples/Solution26.java. Inputs are fixed in the program.  \par Syllabus reading: Reference material. CLO-3.}
\end{frame}

\begin{frame}[fragile,t]{Execution trace}
\small\renewcommand{\arraystretch}{1.5}
\rowcolors{2}{pale}{white}
\begin{tabularx}{\textwidth}{>{\raggedright\arraybackslash}X>{\raggedright\arraybackslash}X>{\raggedright\arraybackslash}X}
\rowcolor{navy}
\color{white}\bfseries Breakpoint & \color{white}\bfseries Observed value & \color{white}\bfseries Interpretation\\
Before loop & maximum = -4 & Valid first candidate\\
After i=1 & maximum = -2 & Larger value replaces it\\
After i=2 & maximum = -2 & -9 cannot replace it\\
Empty input guard & length = 0 & Exception before index access\\
\end{tabularx}
\vspace{3mm}\footnotesize Follow the changing state in execution order.
\note{Manually derived trace for the complete solution. Document that the method requires a nonempty array. Move the maximum calculation into a method returning int. Call it for negative-only data and reject the empty case explicitly.\par Syllabus reading: Reference material. CLO-3.}
\end{frame}

\begin{frame}[fragile,t]{Solution result}
\begin{columns}[T,onlytextwidth]\begin{column}{0.60\textwidth}
-2\par Empty array rejected
\end{column}\begin{column}{0.35\textwidth}\small
Why this result follows

Call it for negative-only data and reject the empty case explicitly.
\end{column}\end{columns}
\note{Expected result: -2
Empty array rejected\par Syllabus reading: Reference material. CLO-3.}
\end{frame}

\begin{frame}[fragile,t]{A common incorrect approach}
int maximum = 0;\par for (int v : values) maximum = Math.max(maximum, v);
\vspace{1mm}\begin{block}{Example}
\begin{lstlisting}
This fails for nonempty arrays containing only negative values. Initialise from the first element under a nonempty-array contract.
\end{lstlisting}
\end{block}
\note{Ask students to predict the failure before discussing the explanation. The right side gives the correction.\par Syllabus reading: Reference material. CLO-3.}
\end{frame}

\begin{frame}[fragile,t]{Boundary and failure checks}
\begin{columns}[T,onlytextwidth]\begin{column}{0.60\textwidth}
Check the stated input assumptions.

Record one debugging session: failing input, hypothesis, observed variable values, correction and regression test.
\end{column}\begin{column}{0.35\textwidth}\small
Reject an empty array with IllegalArgumentException. Start from values[0], traverse from index 1 and return maximum. \{-4,-2,-9\} must give -2.
\end{column}\end{columns}
\note{Use the specification to derive each expected result. Exercise solution: Reject an empty array with IllegalArgumentException. Start from values[0], traverse from index 1 and return maximum. \{-4,-2,-9\} must give -2.\par Syllabus reading: Reference material. CLO-3.}
\end{frame}

\begin{frame}[fragile,t]{Check your understanding}
\begin{columns}[T,onlytextwidth]\begin{column}{0.60\textwidth}
Does rename refactoring intentionally change output? What is a regression test?
\end{column}\begin{column}{0.35\textwidth}\small
Write a prediction and explain the rule that supports it.
\end{column}\end{columns}
\note{Answer: No. It changes a name while preserving behaviour. A regression test checks that a previously corrected failure does not return.\par Syllabus reading: Reference material. CLO-3.}
\end{frame}

\begin{frame}[fragile,t]{Answer and explanation}
\begin{columns}[T,onlytextwidth]\begin{column}{0.60\textwidth}
No. It changes a name while preserving behaviour. A regression test checks that a previously corrected failure does not return.
\end{column}\begin{column}{0.35\textwidth}\small
This fails for nonempty arrays containing only negative values. Initialise from the first element under a nonempty-array contract.
\end{column}\end{columns}
\note{Reconnect the answer to the relevant definition rather than treating it as a fact to memorise.\par Syllabus reading: Reference material. CLO-3.}
\end{frame}


\begin{frame}[fragile,t]{Compare cases and predict outcomes}
\small\renewcommand{\arraystretch}{1.5}\rowcolors{2}{pale}{white}
\begin{tabularx}{\textwidth}{>{\raggedright\arraybackslash}X>{\raggedright\arraybackslash}X>{\raggedright\arraybackslash}X}\rowcolor{navy}
\color{white}\bfseries Observation & \color{white}\bfseries Hypothesis & \color{white}\bfseries Useful inspection\\
Last element omitted & Loop stops too early & Index and length at condition\\
First item skipped & Counter starts at 1 & Initial counter value\\
Total persists & Accumulator not reset & Value before the next run\\
\end{tabularx}\par\vspace{3mm}\begin{exampleblock}{Discuss}Explain each expected result before running the example. Which case would reveal a wrong assumption?\end{exampleblock}
\note{Read the first two columns and invite predictions before discussing the outcome.}
\end{frame}

\begin{frame}[fragile,t]{Apply the idea}
\begin{alertblock}{Independent practice}A sum loop uses i \textless{} values.length - 1 for \{2,4,6\}. Predict the wrong result, identify the missing visit, and correct the condition.\end{alertblock}\vspace{3mm}\begin{enumerate}\item Work through the input by hand.\item Write or correct the program.\item Justify the expected result.\end{enumerate}\note{Allow 3–5 minutes, then compare answers in pairs. A sum loop uses i \textless{} values.length - 1 for \{2,4,6\}. Predict the wrong result, identify the missing visit, and correct the condition.}
\end{frame}

\begin{frame}[fragile,t]{Solution and reasoning}
\begin{lstlisting}
int sum = 0;
for (int i = 0; i < values.length; i++) {
  sum += values[i];
}
\end{lstlisting}\vspace{1mm}\small\begin{itemize}
\item The faulty condition visits indexes 0 and 1 only, giving 6.
\item Index 2, containing 6, is omitted.
\item The corrected loop visits all three elements and gives 12.
\end{itemize}\note{Answer: The faulty condition visits indexes 0 and 1 only, giving 6. Index 2, containing 6, is omitted. The corrected loop visits all three elements and gives 12.}
\end{frame}

\begin{frame}[fragile,t]{Debugging an incorrect array average}
\begin{itemize}
\item The source debugger exercise starts its sum loop at 1 and omits the first element.
\item It also divides ints before returning double, so a fractional mean can be lost.
\item Use a fractional expected result to detect both defects instead of only testing an integer{-}valued mean.
\end{itemize}
\vspace{3mm}\footnotesize\textcolor{accent}{Source: supplied ISB campus slides. See speaker notes and ISB\_Source\_Map.md.}
\note{Adapted from the supplied ISB campus folder: Introduction\&VScode setup.pptx, slides 13. Changed the source five{-}element fixture to four elements so the fractional{-}average defect remains observable; added the empty{-}input check.}
\end{frame}

\begin{frame}[fragile,t]{Worked problem: Debugging an incorrect array average}
\begin{block}{Task}Correct the sum and average for \{1,2,3,4\}. Inspect i, sum and the final divisor with a breakpoint in the loop.\end{block}
\small Input: Values are fixed in the program for a reproducible demonstration.\par\vspace{3mm}\begin{exampleblock}{Before running}Predict the output and identify the variables that change.\end{exampleblock}
\note{Adapted from the supplied ISB campus folder: Introduction\&VScode setup.pptx, slides 13. Changed the source five{-}element fixture to four elements so the fractional{-}average defect remains observable; added the empty{-}input check.}
\end{frame}

\begin{frame}[fragile,t]{Complete ISB example (1/1)}
\begin{lstlisting}
public class IsbLecture26 {
  public static void main(String[] args) throws Exception {
    int[] values = {1, 2, 3, 4};
    if (values.length == 0) throw new IllegalArgumentException("empty");
    int sum = 0;
    for (int i = 0; i < values.length; i++) sum += values[i];
    double average = (double) sum / values.length;
    System.out.println(sum);
    System.out.println(average);
  }

}
\end{lstlisting}
\note{Adapted from the supplied ISB campus folder: Introduction\&VScode setup.pptx, slides 13. Changed the source five{-}element fixture to four elements so the fractional{-}average defect remains observable; added the empty{-}input check. Complete file: ISB\_Java\_Examples/IsbLecture26.java.}
\end{frame}

\begin{frame}[fragile,t]{Worked trace and expected result}
\small\renewcommand{\arraystretch}{1.4}\rowcolors{2}{pale}{white}
\begin{tabularx}{\textwidth}{>{\raggedright\arraybackslash}X>{\raggedright\arraybackslash}X>{\raggedright\arraybackslash}X}\rowcolor{navy}
\color{white}\bfseries Implementation & \color{white}\bfseries Sum & \color{white}\bfseries Returned average\\
Starts at 1; integer division & 9 & 2.0\\
Index fixed only & 10 & 2.0\\
Index and division fixed & 10 & 2.5\\
\end{tabularx}\par\vspace{2mm}
\begin{block}{Expected console output}\ttfamily\footnotesize 10\par 2.5\end{block}
\note{Adapted from the supplied ISB campus folder: Introduction\&VScode setup.pptx, slides 13. Changed the source five{-}element fixture to four elements so the fractional{-}average defect remains observable; added the empty{-}input check. Trace and expected values independently worked for this edition.}
\end{frame}

\begin{frame}[fragile,t]{Extend the ISB example}
\begin{alertblock}{Practice}Why might \{1,2,3,4,5\} hide the integer{-}division defect after the loop is fixed?\end{alertblock}\vspace{3mm}\begin{exampleblock}{Explain your reasoning}Its exact mean is 3.0, so integer division happens to produce the same numeric answer.\end{exampleblock}
\note{Adapted from the supplied ISB campus folder: Introduction\&VScode setup.pptx, slides 13. Changed the source five{-}element fixture to four elements so the fractional{-}average defect remains observable; added the empty{-}input check. Answer: Its exact mean is 3.0, so integer division happens to produce the same numeric answer.}
\end{frame}
\begin{frame}[fragile,t]{Lecture summary}
\begin{columns}[T,onlytextwidth]\begin{column}{0.60\textwidth}
\begin{itemize}
\item breakpoints and variable inspection
\item debugging from refactoring
\item Refactor while preserving observable behaviour
\end{itemize}
\end{column}\begin{column}{0.35\textwidth}\small
\begin{itemize}
\item Document that the method requires a nonempty array.
\item Move the maximum calculation into a method returning int.
\item Call it for negative-only data and reject the empty case explicitly.
\end{itemize}
\end{column}\end{columns}
\note{Ask students to give one concrete example for the concepts listed. Use the worked solution to connect the topics.\par Syllabus reading: Reference material. CLO-3.}
\end{frame}

\begin{frame}[fragile,t]{After-class practice}
\begin{columns}[T,onlytextwidth]\begin{column}{0.60\textwidth}
Record one debugging session: failing input, hypothesis, observed variable values, correction and regression test.
\end{column}\begin{column}{0.35\textwidth}\small
Syllabus reading\par Reference material

Next lecture: Library components and APIs
\end{column}\end{columns}
\note{Reading labels follow the supplied outline. Check topic headings against the available textbook edition. Require a program and expected results for the practice task.\par Syllabus reading: Reference material. CLO-3.}
\end{frame}
\end{document}
